\documentclass[12pt]{article}
\usepackage[utf8]{inputenc}
\usepackage[T1]{fontenc}
\usepackage[margin=1in]{geometry}
\usepackage{setspace}
\usepackage{xcolor}
\usepackage[colorlinks=true, urlcolor=blue!60!black, linkcolor=blue!60!black]{hyperref}
\setlength{\parindent}{2em}
\setlength{\parskip}{1.25em}
\renewcommand{\baselinestretch}{1.0}
\begin{document}
\begin{center}
{\LARGE \textit{West Virginia v. B. P. J. / Little v. Hecox}, Nos. 24-43 \& 24-38 (2026)\par}
\vspace{0.8em}
{\large Case Brief --- Four-Page Edition\par}
\end{center}
\vspace{1em}

\section{Memory Jogger}

In consolidated cases decided June 30, 2026, the Court held, 9--0 in judgment under Title IX and 6--3 under the Equal Protection Clause, that States may restrict public-school women's and girls' sports teams to biological females without individualized exceptions for transgender girls and women (\href{https://www.law.cornell.edu/supremecourt/text/24-43}{Cornell LII}).

\section{Detailed Case Facts}

Between 2020 and 2026, 27 States restricted women's and girls' school teams to biological females, and the NCAA, USOPC, and IOC adopted parallel policies. West Virginia's Save Women's Sports Act (2021) designated teams by biological sex'' and barred biological males from female teams in contact sports or sports selected by competitive skill. B. P. J. (Becky Pepper-Jackson), assigned male at birth, socially transitioned in third grade, took puberty blockers to prevent male puberty, and began estrogen in sixth grade. Barred from girls' cross-country and track, she competed under court injunctions and later won a state championship in girls' shot put. Idaho's Fairness in Women's Sports Act (2020) similarly barred biological males from female teams from primary school through college. Lindsay Hecox, a transgender woman who underwent male puberty, sought to try out for Boise State's women's track and cross-country teams.

\section{Procedural History}

In West Virginia, the district court first enjoined the Act, then granted summary judgment to the State; the Fourth Circuit reversed, directing judgment for B. P. J. on Title IX and remanding the equal protection claim for factfinding on athletic advantage. In Idaho, the district court preliminarily enjoined the Act and the Ninth Circuit affirmed. The Supreme Court granted certiorari in July 2025, rejected Hecox's later suggestion of mootness, heard argument on January 13, 2026, and reversed both judgments.

\section{Judicial Votes}

Reversed and remanded: \textbf{9--0} in judgment on Title IX; \textbf{6--3} on Equal Protection and mootness.
\begin{itemize}
\item \textbf{Majority:} Justice Kavanaugh, joined by Chief Justice Roberts and Justices Thomas, Alito, Gorsuch, and Barrett.
\item \textbf{Concurrences:} Justices Thomas and Gorsuch, each joining in full.
\item \textbf{Concurring in the judgment in part and dissenting in part:} Justice Sotomayor, joined by Justices Kagan and Jackson; Justice Jackson separately.
\end{itemize}

\section{Holding}

(1) Under Title IX, the 1974 Javits Amendment, and the 1975 regulations, sex'' means biological sex, and schools may keep women's teams limited to biological females. (2) Such state laws classify by biological sex, not transgender status, and satisfy intermediate scrutiny as substantially related to safety and competitive fairness; States need not grant individualized or subclass exemptions. \textit{Hecox} was not moot. The Court reserved whether States may \textit{include} transgender girls on female teams, whether transgender status triggers heightened scrutiny, and how \textit{Bostock} applies outside athletics.

\section{Analysis of Opinions}

\textbf{Majority (Kavanaugh).} Enduring'' physical differences in height, strength, speed, and endurance justify sex-separated teams. The ordinary meaning of sex'' in 1972, confirmed by the contemporaneous regulations, is biological; because sports are generally zero sum,'' a biological line is reasonable.'' Title VII and \textit{Bostock} differ because Title IX expressly authorizes separate teams. On equal protection, under \textit{Nguyen} and \textit{Michael M.}, the plaintiffs' as-applied challenge was a backdoor argument for strict scrutiny'' requiring unadministrable athlete-by-athlete comparisons; citing \textit{Skrmetti}, the laws survive any level of scrutiny amid ongoing medical and scientific debate.''

\textbf{Thomas (concurring).} Transgender status is not a suspect class; [m]en and boys with gender dysphoria are not women or girls, even if they believe that they are.''

\textbf{Gorsuch (concurring).} Spending Clause legislation must give clear notice under \textit{Pennhurst}, and \textit{Bostock} is consistent because separate teams are authorized, not unlawful discrimination.

\textbf{Sotomayor (in part dissenting).} B. P. J.'s Title IX claim failed only because of her concession, and \textit{Bostock} still protects transgender students elsewhere; \textit{Hecox} was moot. The Court should have allowed factfinding on whether a discrete subclass, transgender girls who never underwent male puberty, furthers the State's interests, and it collapsed intermediate scrutiny into rational-basis deference.

\textbf{Jackson (in part dissenting).} Sex-stereotyping precedents show that Title IX protects gender identity, not only sex assigned at birth.

\section{Examples: Future Applications}

\begin{enumerate}
\item \textbf{Same-Side (Upheld): Intramural Basketball.} A State limits public university women's intramural and club basketball to biological females. A categorical line is permitted.
\item \textbf{Same-Side (Upheld): Girls' Swimming.} A state athletic association applies its rule based on biological sex to a transgender girl who took puberty blockers. \textit{B. P. J.} forecloses the as-applied challenge.
\item \textbf{Opposite-Side (Invalid): Coed Chess or Debate.} A district bars a transgender student from coed activities. That falls outside Title IX's athletic carve-out.
\item \textbf{Opposite-Side (Invalid): Expulsion.} A school expels a student for socially transitioning. Under \textit{Bostock}'s logic, that is discrimination on the basis of sex.''
\item \textbf{Fence-Sitter: An Inclusive State Law.} Cisgender athletes challenge a State's guarantee that transgender girls may play on girls' teams. Footnote 2 left that question open.
\end{enumerate}

\section{Critique}

\textbf{Originalist and textualist praise.} Commentators praise the Court for adhering to the 1972 meaning of sex,'' enforcing clear-notice limits, and respecting democratic federalism, while refusing to make courts national athletic commissions supervising hormone-based exemptions.

\textbf{Progressive critique.} Erwin Chemerinsky warns that by omitting \textit{Virginia}'s exceedingly persuasive justification'' standard and rejecting individualized fit, the majority watered down intermediate scrutiny for all sex discrimination. Craig Konnoth argues that plaintiffs litigated the wrong issue by conceding biological sex rather than challenging whether statutory definitions track performance. The \textit{Harvard Law Review} faults all five opinions for ignoring the coordinated campaign behind the bills.

\section{Key Quotations}

\begin{quote}
``Consistent with Title IX and the Equal Protection Clause, we hold that the States may maintain women's and girls' sports for biological females. They may determine eligibility for women's and girls' sports based on biological sex. The Constitution and Title IX do not require an overhaul of women's and girls' sports throughout America.''
\end{quote}
\noindent --- Justice Kavanaugh, for the Court, slip op. at 29.

\begin{quote}
``The underlying medical and scientific premise of the plaintiffs' entire equal protection argument is that at least some biological males who identify as female and take puberty blockers or hormones do not retain physical advantages over biological females. Even if true, that empirical claim would not alter the equal protection conclusion set forth above.''
\end{quote}
\noindent --- Justice Kavanaugh, for the Court, slip op. at 24.

\begin{quote}
``Because of the Court's decision today, West Virginia, and any other state actor, can deny B. P. J. and others like her these experiences simply because it thinks they have an inherent athletic advantage, even if the facts show that they do not. In the end, to the Court, the facts do not matter, even though the consequences are serious.''
\end{quote}
\noindent --- Justice Sotomayor, concurring in the judgment in part and dissenting in part.

\end{document}
